\begin{textsample}{NEG Dim 6 – global\_south\_2025\_02 – Score -1.00 – global\_south\_2025\_02/120967219.txt}  \label{ex:f6_neg_036}
THE strangeness of the idea can not be denied.
United States President Mr.
Donald Trump has said, in effect, that the US may take over the Gaza Strip and develop the area since the current fragile peace may not hold for long.
No matter the words of promise of development uttered by Mr.
Trump, the idea is nothing less than utopian since Mr.
Trump is thinking of"owning"the war-ravaged Gaza Strip and introduce a strong \textbf{dose} of economic development so that good numbers of jobs are available to the people there and peace would be the essential condition for that reign to continue.
He might not have said all this in these many words, all right, but there is no doubt that Mr.
Trump ' s idea has the potential to shake the world at large owing to the implication -- and complication -- everybody may sense and see differently.
Of course, in his second coming, Mr.
Donald Trump is showering upon the world shocks after shocks with his bold initiatives that he had promised @ @ @ @ @ @ @ @ @ @ of owning the Gaza Strip, Mr.
Trump appears to have been more than candid in his discussions with Israeli Prime Minister Mr.
Benjamin Netanyahu in Washington DC.
The Israeli leader might not have agreed with the US President word by word, but he -- or anybody -- would have only lame reasons to oppose the idea of economic development that would act as a stronger shield for peace to keep prevailing.
But as the world would know what Mr.
Donald Trump may mean ultimately, everybody would be alarmed -- to say the least -- by the dare-devilry involved in such a thought.
Only Mr.
Trump appears capable of bringing to fore such ideas.
We may offer another description for this latest from Mr.
Trump.
In the first flush, we have called the idea"strange".
In second thought, we may even describe it as quixotic ( in deference to author Miguel de Cervantes who imagined the character of Don Quixote -- in the 17th century -- who has his own way of looking at the world and @ @ @ @ @ @ @ @ @ @ air, so to say ).
We choose to use that metaphor for Mr.
Trump because we do -- as does the world -- realise that he can unleash his imagination in strangest possible ways.
There should actually be no doubt that if the US tries to do something as per Mr.
Trump ' s idea in the Gaza Strip, the action would trigger a chain reaction whose speed and trajectory may not remain in anybody ' s control.
It is clear that by allowing such thoughts to surface, Mr.
Trump has given the collective global leadership a good reasons to scratch heads and look askance at one another.
Who knows, any move that direction may trigger a conflict far more serious and dangerous in nature than the current one.
But, if at all the move happens to go on stream, even in a half-baked manner, then the nature and colour of global international realpolitik would undergo a dramatic and even drastic change.
In still other words, the world may find itself close to the brink of another @ @ @ @ @ @ @ @ @ @ never before.
The seeds of this apprehension lie in Mr.
Trump ' s reported statement that the US may send troops to Gaza Strip if"necessary".
He also said during that statement that currently the Gaza Strip is like a"hell hole"-- which needs to be corrected"not for Israel but ( also ) for the entire Middle East".
Anything can be read out of such statements more particularly when someone of Mr.
Trump ' s style and signature is saying those.
No matter the different scenarios that can be worked out of such statements, there is no doubt that the idea has some substance outside the realm of global geopolitics -- if the world has the inclination to think deeply about it.

% matched lemmas: dose
\end{textsample}
